\chapter{Literature Review Summary Table}

\begin{longtable}{|p{3.2cm}|p{3cm}|p{4.8cm}|p{4.8cm}|}
\caption{Summary of prior studies on multimodal and context-aware misinformation detection}
\label{tab:literature_summary} \\
\hline
\textbf{Study} & \textbf{Focus} & \textbf{Key finding / contribution} & \textbf{Main limitation} \\
\hline
\endfirsthead

\hline
\textbf{Study} & \textbf{Focus} & \textbf{Key finding / contribution} & \textbf{Main limitation} \\
\hline
\endhead

\hline
\endfoot

\hline
\endlastfoot

Wang et al. (2018) -- EANN
& Early multimodal fake news detection
& Showed that combining text + image features improves misinformation detection over text-only methods
& Fusion was still relatively early-stage and did not fully model deeper cross-modal relationships \\
\hline

Nakamura et al. (2020) -- Fakeddit
& Multimodal benchmark dataset
& Provided a large benchmark showing that multimodal classifiers outperform unimodal baselines
& Focused more on benchmarking than integrating social context \\
\hline

Xue et al. (2021) -- MCNN
& Text-image consistency modelling
& Demonstrated that modelling the relationship between text and image improves classification
& Mainly content-focused, with limited use of propagation or user context \\
\hline

Qian et al. (2021) -- HMCAN
& Hierarchical multimodal attention
& Showed that attention-based multimodal fusion captures cross-modal interactions better than simple feature concatenation
& Still concentrated on content-level signals rather than broader platform context \\
\hline

Song et al. (2021)
& Social/propagation context
& Found that diffusion patterns and user interaction signals are useful for misinformation detection
& Reliant on interaction data, which may be sparse in early detection \\
\hline

Min et al. (2022)
& Temporal and behavioural signals
& Confirmed that temporal dynamics and engagement behaviour can improve classification
& Less effective when user activity data is limited or delayed \\
\hline

Segura-Bedmar \& Alonso-Bartolome (2022)
& Multimodal benchmarking
& Reinforced evidence that multimodal approaches outperform text-only models on social media misinformation tasks
& Does not resolve how multimodality compares when social context is held constant \\
\hline

Nadeem et al. (2023) -- EFND
& Multimodal + contextual integration
& Showed that combining content and contextual signals can produce strong detection performance
& Hard to isolate whether gains come from multimodality or the added context \\
\hline

Jing et al. (2023) -- MPFN
& Progressive multimodal fusion
& Demonstrated that better fusion design, not just adding more modalities, improves results
& Limited social-context integration \\
\hline

Truică et al. (2024) -- DANES
& Context-aware ensemble model
& Combined a text-based content branch with a social/propagation branch, achieving strong benchmark performance
& Content branch remains text-only, so visual misinformation may be missed \\
\hline

Wang et al. (2024)
& Joint multimodal-context modelling
& Supported the view that multimodal content and social context are complementary
& Social context representation differs from other studies, making comparison difficult \\
\hline

Zhu et al. (2024)
& Propagation-aware detection
& Showed that behavioural and diffusion evidence adds strong classification value beyond content alone
& Performance depends on availability of network/interaction data \\
\hline

Truică et al. (2025) -- GETAE
& Context-aware ensemble model
& Similar to DANES, achieved strong results by combining text content + social context
& Still inherits the weakness of a text-only content branch \\
\hline

Jing et al. (2025)
& Unified multimodal-context model
& Reported strong performance from joint modelling of text, visuals, and context
& Multiple components change at once, so the exact source of improvement is unclear \\
\hline

Chen et al. (2025)
& Behaviour/context-enhanced multimodal model
& Found that integrating behavioural context with multimodal content improves detection
& Comparison across studies is difficult because context features vary widely \\
\hline

Sun et al. (2025)
& Dynamic graph/context fusion
& Showed value in modelling dynamic relationships and graph-based context alongside content
& Adds complexity and reduces comparability with simpler baselines \\
\hline

Dellys et al. (2025)
& Comparative/contextual discussion
& Highlighted the difficulty of identifying the specific contribution of multimodal substitution in context-aware systems
& Does not fully isolate multimodal effects within DANES/GETAE-style architectures \\
\hline

Ji et al. (2026)
& Recent integrated framework
& Further supported the effectiveness of combining multimodal and contextual signals
& The separate effect of multimodality is still not clearly decomposed \\
\hline

\end{longtable}